\documentclass[12pt]{article}
\usepackage{graphicx} % Required for inserting images
\usepackage[portuguese]{babel}
\usepackage{url}
\usepackage[hidelinks]{hyperref}
\usepackage{xurl} % must be loaded AFTER hyperref
\usepackage{csquotes}
\usepackage[backend=biber, style=abnt]{biblatex}
\usepackage{bookmark}
\usepackage{epigraph}
\addbibresource{Artigo_ESG.bib}
\setlength{\epigraphwidth}{0.5\textwidth}


%\usepackage[alf]{abntex2cite}
%\bibliographystyle{abntex2-alf}

\title{A esquizofrenia do ESG: entre a ética do capital e a ética da vida (\textbf{ESBOÇO})}
\author{Nicholas Funari Voltani}
\date{19 de fevereiro de 2026}

\begin{document}
\maketitle

%\epigraph{``...\textit{for certainly, you were better take for business a man somewhat absurd than over-formal.}'' (Francis Bacon)}




\section{Introdução}
A crise climática que aflige o mundo no século XXI naturalmente traz o desejo de, ou o desespero por, formas mais sustentáveis de convívio do homem com a natureza, dentre as quais o combate ao desmatamento, mitigação de emissões de gases de efeito estufa, a busca por uma economia \textit{net-zero} de carbono\footnote{Ou seja, em que todas as emissões de carbono são contrabalançadas por pela capacidade do planeta Terra de absorvê-las.}, etc. Ademais, a questão climática naturalmente tem rastros sociais: populações marginalizadas econômica/politicamente são as mais afetadas pelas consequências climáticas. 

\textbf{Menção dos tipos de ``Sustainable Investing'' em \cite{Fichtner2024} à guisa de introdução mínima ao ESG. Menção da bagunça e elusiveness do conceito. Algo a mais...} 



Mas, ao mesmo tempo, a sociedade capitalista pressupõe, sub-repticiamente, uma ontologia conservadora sobre a conformação do mundo em que habita, qual seja, de que o mundo nada mais é do que aquilo que pode ser observado --- e tudo aquilo que não pode ser observado, simplesmente não pode \textit{ser} ---, o que Bhaskar chama de “realismo empírico” \cite{Bhaskar1978}. Ou seja, a ontologia pressuposta pela sociedade capitalista ignora tanto fenômenos efetivos mas não-observáveis, quanto os mecanismos geradores dos fenômenos efetivos, escanteando-os como “metafísicos”, meramente “ideais”.

As consequências lógicas de uma ontologia empirista decorrem de seu enclausuramento em um mundo não-mais-que-empírico --- portanto, de um mundo ``\textit{fechado}'' --- são que a ciência torna-se algo meramente instrumental, e que os fenômenos sociais não têm como ser algo a mais do que as ações individuais.

É neste contexto todo que elucida-se melhor onde o ESG insere-se juntamente às contradições do mundo capitalista, e como seus anseios (porventura genuínos) de mudança sustentável são afogados por sua alma indelevelmente capitalista.
\textbf{Fazer alguma abordagem mínima sobre ética e moral no que tange ao capital, pôr teleológico... Não necessariamente aqui na introdução, cabe mais no meio do texto, vide abaixo. \cite{Medeiros2013,Medeiros2013a,Medeiros2016}} 






\section{A doutrina Friedman}
Em seu famoso artigo de 1970 pelo New York Times, Milton Friedman faz seu argumento contra as ``responsabilidades sociais'' que começavam a ser imputadas ao homem de negócios\footnote{``Na verdade, eles [os homens de negócios ``socialmente comprometidos''] estão --- ou estariam, se eles ou qualquer outra pessoa os levasse a sério --- pregando o socialismo puro e simples [!]. Empresários que falam dessa maneira são marionetes involuntárias das forças intelectuais que vêm minando os alicerces de uma sociedade livre nas últimas décadas.'' \cite{Friedman1970}.}, estabelecendo, dessa forma, a assim chamada `` doutrina Friedman''. Encapsulada na última frase de sua publicação, ele afirma que, em uma ``sociedade livre'',
\begin{quote}
    ``há uma e somente uma responsabilidade social dos negócios --- usar os recursos destes [negócios] e engajar-se em atividades destinadas a aumentar seus lucros [dos negócios], desde que estejam \textit{dentro das regras do jogo}, ou seja, engajar-se em competição aberta e livre, sem enganação ou fraude.'' \cite[grifo meu]{Friedman1970}
\end{quote}


O eixo central de sua argumentação é que um executivo corporativo, tendo sido apontado pelos acionistas (\textit{shareholders}) de uma empresa, torna-se responsável pela gestão e administração dos ativos destes, e, naturalmente, conduzirá os negócios de acordo com \textit{os interesses destes}. Dessa forma, caso deseje engajar-se em ``responsabilidades sociais'' de forma que elas contraponham-se aos interesses de seus acionistas --- quais sejam, essencialmente a geração de lucros ---, então ele está ``gastando o dinheiro deles'' --- ou, mais explicitamente, está desperdiçando-o, está agindo de má fé para com aqueles que apontaram-no para esta responsabilidade. Se ele deseja engajar-se socialmente, que o faça em sua vida privada, com seus próprios recursos. Afinal de contas, a sociedade ``não passa de uma soma de indivíduos'', e a mudança ``só pode vir dos indivíduos''.\footnote{``As discussões sobre as “responsabilidades sociais das empresas” são notáveis pela sua imprecisão analítica e falta de rigor. O que significa dizer que “empresas” têm responsabilidades? \textit{Somente pessoas podem ter responsabilidades}. Uma corporação é uma pessoa jurídica e, nesse sentido, pode ter responsabilidades artificiais, mas não se pode dizer que “empresas” como um todo tenham responsabilidades, nem mesmo nesse sentido vago. O primeiro passo para a clareza na análise da doutrina da responsabilidade social das empresas é perguntar precisamente o que ela implica e para quem.'' \cite[grifo meu]{Friedman1970}.}





\section{O trabalho corporativo e sua ética}

\begin{quote}
    ``...mesmo o próprio trabalho comercial que, à primeira vista, parece desprovido de atração [além do ``ganho'' financeiro], dá muitas vezes \textit{um grande prazer oferecendo um objetivo ao exercício das faculdades do homem e a seus instintos de emulação e de poder}[,] pois, assim como um cavalo de corrida ou um atleta, que exige tudo de cada um dos seus nervos para exceder seus concorrentes, e sente prazer nesse esforço, assim também um industrial ou um comerciante é muitas vezes estimulado mais \textit{pela esperança de vencer seus rivais} do que pelo desejo de juntar algo à sua fortuna.'' \cite[p. 39; grifo meu]{Marshall1985}
\end{quote}

Os altos salários dos trabalhadores corporativos já dão boa conta de suas necessidades materiais básicas, dando-lhes margem para buscar novas satisfações neste trabalho que, em sua essência, não lhe diz respeito\footnote{``Além do esforço dos órgãos que trabalham, a atividade laboral exige a vontade orientada a um fim, que se manifesta como atenção do trabalhador durante a realização de sua tarefa, e isso tanto mais quanto menos esse trabalho, pelo seu próprio conteúdo e pelo modo de sua execução, atrai o trabalhador, portanto, quanto menos este último usufrui dele como jogo de suas próprias forças físicas e mentais.'' \cite[p. 256]{Marx2017}}. Conforme adequa-se ao \textit{habitus} do meio em que trabalha, começa a perceber, mesmo que tacitamente, suas regras e seu ``espírito''. É neste sentido que mesmo os trabalhadores assalariados corporativos em cargos mais ``baixos''\footnote{Desconsiderando cargos mais, entre aspas, ``braçais''. Qualquer um que trabalhou na Faria Lima sabe que os garçons de restaurantes e atendentes de balcão tendem a ser mulheres negras (ou então, homens negros). O completo oposto é o que se observa nos cargos de ``colarinho branco''. É quase cruel dizer que há uma ``igualdade de oportunidades'' entre aqueles que servem a mesa num restaurante do Itaim Bibi e aqueles em que nele almoçam.}, ínfimas peças deste sistema que eles já encontram ``pronto'', tornam-se capazes de personificar a concorrência capitalista. O que os move, satisfeitas suas necessidades e preocupações com sua própria sobrevivência, pode agora ser meramente o \textit{thrill of the hunt}, a adequação de agir conforme as ``regras do jogo'' do mercado, de vencer pela astúcia, seja ela \textit{fair play} ou um \textit{foul play} aceitável (seja legalmente ou dentro do orçamento jurídico). 

Os salários dos neófitos deste meio têm, também, de ser altos o suficiente, para que possam ser sua porta de entrada neste mundo que lhes é alheio; no outro lado da moeda, os salários dos gerentes têm de ser altos o suficiente para que não tenham razão de procurar a porta de saída da empresa, i.e. para que não sintam-se ``alheios'' na empresa onde estão. Para os primeiros, tais salários são necessários; para os últimos, são suficientes. Os primeiros agem em nome da empresa, mas em prol de si mesmos; os últimos agem por si mesmos --- por sua própria satisfação de ``vencer seus competidores'' ---, mas em prol da empresa\footnote{Embora ele até possa esquecer-se disso, iludindo-se com a própria máscara que pôs no rosto, ele é eventualmente forçado a aceitar essa egoísta realidade quando ``sua'' empresa desfaz-se dele, análogo a como ele desfez-se de sua própria vida pessoal por lealdade à empresa. O capital é prudente o suficiente de não apegar-se demasiado a nenhuma das ferramentas com as quais compraz-se, sejam elas de carne ou de metal.}.

Dessa forma, o trabalhador corporativo, de alto ou baixo escalão, é capaz de incorporar a ética do mercado tão aderentemente que até torna-se capaz de ``entrar em fluxo'' com ele, de \textit{entendê-lo} e sintonizar-se às ações que são \textit{adequadas} a se tomar. Alguns destes trabalhadores podem até mesmo chegar a crer que suas ações, por serem tão paralelas ao que espera-se do mercado, moldam-no \textit{diretamente}; chegam a entreter a fantasia de que é o indivíduo que molda a sociedade de punho próprio, e que esta não consiste em mais do que a soma de seus espécimes e de suas ações --- embora saibam, bem no fundo, irrefletidamente, que falar ``do mercado'' é algo além de falar ``de seus demais concorrentes'' ou da soma dos \textit{players} de um setor da indústria. Isso se manifesta quando falam, por exemplo, de \textit{ciclos de negócios} (\textit{business cycles}), que \textit{impõem-se às empresas} e \textit{forçam sua adaptação}, não o contrário!\footnote{Lukács coloca bem como o indivíduo corporativo torna-se subjetivamente afetado pela reificação do capitalismo e de suas ``leis de funcionamento'': ``O comportamento do homem esgota-se, portanto, no cálculo correto das oportunidades desse curso [\textit{Ablauf}: curso, processo] (cujas ``leis'' ele já encontra ``prontas''), na habilidade de evitar os ``acasos'' perturbadores [do \textit{business as usual}] por meio da aplicação de dispositivos de proteção e medidas defensivas (que se baseiam igualmente na consciência e na aplicação de ``leis'' semelhantes); muitas vezes, [o homem] chega até mesmo a se deter no cálculo das probabilidades dos possíveis efeitos de tais ``leis'', sem sequer tentar intervir no próprio processo pela aplicação de outras ``leis'' (como nos esquemas de segurança etc.). Quanto mais se considera essa situação em profundidade e independentemente das lendas burguesas sobre o caráter ``criador'' dos expoentes da época capitalista, tanto mais claramente aparece, em tal comportamento, a analogia estrutural com o comportamento do operário em relação à máquina que ele serve e observa, e cujo funcionamento ele controla enquanto a contempla.'' \cite[p. 218-9; destaques meus]{Lukacs2003a}.}







\section{A ética do capital}
A ética do capital tem por \textit{télos} sua autovalorização, priorizando, portanto, os modos de ação que conduzam-no à sua \textit{raison d'être}. É neste sentido que essas falas acima mencionadas, de tais marechais do exército do capital, expõem a hipocrisia do ESG no meio corporativo-financeiro: não como uma denúncia de sua negligência para com a sociedade, e sim de sua negligência para com os negócios; não porque suas ações --- que restringem-se a precificações e negócios financeiros --- são só ``para inglês ver'', e sim porque são uma perda de tempo, recursos e pessoal \textit{no que tange ao próprio capital}. 

Assim como as leis burguesas não se exaurem nos livros que as contam, nem nas \textit{fake news} que as deturpam, a ética do capital não exaure-se na ética de seus atores, nem em seus discursos contraditórios às suas crenças e ações.\footnote{``Quanto às pessoas artificiais, em certos casos algumas das suas palavras e ações pertencem àqueles a quem representam.

Nesses casos a pessoa é o ator, e aquele a quem pertencem as suas palavras e ações é o AUTOR, casos estes em que o ator atua por autoridade.'' (HOBBES, p. 138, Cap. XVI). \textbf{Cabe checar, pode ser uma ótima analogia.}} 

Portanto, não criticam o ESG porque a alma deste paradigma é antissocial, e sim porque ela é anti-capital. Como conciliar as demandas por mudança social se a disrupção da sociedade vem da própria opressão do capital sobre a sociedade? Pela sua hegemonia, o capital só cumprirá sua ``responsabilidade social'' por ``boa vontade'', o que é o mesmo que dizer ``por motivos ulteriores'', pois sabe que age contra seu espírito ao fazê-lo; sabe que peca perante seu Deus-dinheiro e que precisa seguir Seu evangelho até o fim dos tempos.

A própria concorrência impele os capitais individuais a não se preocuparem com o futuro tão-tão-distante de alguns anos adiante\footnote{``Nosso prazo médio de empréstimo [na HSBC] é 6 anos. O que acontece com o planeta no 7º ano é irrelevante para nossa carteira de crédito.'' (Kirk, \textit{op. cit.}, 11:05). ``Quem hoje está investindo e analisando prefere não se preocupar porque não sabe se estará daqui a algumas décadas na mesma área ou sequer vivo[;] a pessoa dá o check favorável para a empresa e pronto. (sic)'', diz Fabio Alperowitch em \cite{Angelo2022}.}, pois são os capitais de hoje com que disputa-se no Valhalla que é o mercado. [em que lutam as mais aguerridas batalhas de dia, e sentam-se à mesa juntos, para o jantar, após o fechamento da Bolsa.]

\textbf{Adicionar anotações de Lukács} 








\section{Ambiental, Social, Governança: o paradigma ESG/ASG}
Sobre Stuart Kirk:\footnote{É com este causo que Telésforo introduz seu texto em \cite{Telesforo2025}.}
\begin{itemize}
    \item Discurso no evento do Financial Times, FT Moral Money Summit \cite{Kirk2022} (20 de maio de 2022)
    \item O próprio CEO do group HSBC, Noel Quinn, posta no seu Linkedin --- num sábado! --- sobre seu rechaço ao discurso de Kirk no dia anterior \url{https://www.linkedin.com/posts/noel-quinn-hsbc_i-do-not-agree-at-all-with-the-remarks-activity-6933729465193668608-14ho?utm_source=social_share_send&utm_medium=member_desktop_web&rcm=ACoAADGDblYBTpg3cdnCuk-ijUGghBmcJcxAlLc}
    \item Kirk é suspenso da HSBC em 22 de maio de 2022 \cite{Makortoff2022}
    \item Kirk se demite da HSBC em 7 de julho de 2022 \cite{Reuters2022}
    \item Kirk é contratado em novembro de 2022... pelo próprio \textit{Financial Times}! \cite{FinancialTimes2022}
\end{itemize}


Ora, se Kirk realmente fosse considerado um desvairado por seus pares, ele não teria sido contratado pelo próprio jornal --- respeitado na área de Economia e finanças --- em cujo evento ele ``disse barbaridades''! 

E --- que surpresa! --- o HSBC já recuou em seus compromissos ESG, frente à pressão, abre aspas, ``anti-\textit{woke}'' que os EUA assumiram após a posse de Trump em 2025: protelou em 20 anos seu louvável objetivo de alcançar o mágico \textit{net-zero}, de 2030 para 2050 \cite{Segal2025a}, e saiu do grupo \textit{Net-Zero Banking Alliance} \cite{Segal2025}. Tudo isso no contexto do segundo mandato de Trump e sua cruzada contra qualquer coisa relacionada com ``inclusão'', ``sustentabilidade'' etc. Isso já é relatado em \cite{Mirza2025}, dez dias após ação presidencial da Casa Branca para, abre aspas, ``Acabar com programas/preferências radicais e onerosos de DEI\footnote{Diversidade, Igualdade e Inclusão (\textit{Diversity, Equity and Inclusion}).} [\textit{Ending Radical And Wasteful Government DEI Programs And Preferencing}]'' \cite{WhiteHouse2025}.

Uma coisa é cortar a visibilidade destes programas para o público, mas continuá-los internamente. Outra é ver isso como uma oportunidade de desfazer-se de fardos desnecessários e expandir seus negócios sem peso na consciência. 



%Tomemos um exemplo: suponhamos que o CEO de algum banco tenha sido influenciado pelo temível ``esquerdismo'' e descubra as atrocidades cometidas em nome do agronegócio, decidindo que não quer mais financiar suas atrozes atividades. Suponhamos então que nosso honrável executivo formaliza algum plano concreto que corte todos seus negócios com o agronegócio (e suponhamos que ``agronegócio'' seja bem delimitado dentro do portfolio do banco). Naturalmente, além do ultraje que seus clientes do agronegócio sentiriam, também alguns de seus \textit{outros} clientes poderiam tornar-se temerosos: e se \textit{minha} área de atividade for a próxima a ser cortada? Resultado: provável fuga indireta de clientes, \textit{bank runs}. E além disso, e talvez mais importante, a \textit{imensa} depreciação do valor de mercado deste banco. 

Pode até ser que ``tudo valha no amor e na guerra'', mas isso certamente não vale para os negócios. Não é preciso ler o artigo de Friedman para saber que um executivo que siga suas predileções pessoas em detrimento das intenções de seus acionistas, com o dinheiro de seus acionistas, teria seus dias contados em sua empresa (quem dirá em sua carreira), mas, como Friedman colocou-o tão candidamente em sua publicação, vale ler dele próprio:
\begin{quote}
    ``E, mesmo que ele [o executivo corporativo] queira, ele conseguiria mesmo se safar de gastar o dinheiro de seus acionistas, clientes ou empregados? Não iriam esses acionistas demiti-lo? (Seja os [acionistas] atuais ou aqueles que assumirão depois que suas [ó, tão nobres!] ações em nome da responsabilidade social tenham reduzido os lucros da corporação e o preço de suas ações.) Seus clientes e empregados o deserdariam por outros produtores e empregadores menos escrupulosos [!] em exercer suas responsabilidades sociais [\textit{their social responsabilities}].'' \cite[destaques meus]{Friedman1970}
\end{quote}

Isso até mesmo os homens de negócios do século XVII já sabiam, como bem resumido por Francis Bacon: ``É certamente melhor convidar para os negócios um homem algo obtuso [\textit{somewhat absurd}] do que um que seja formal demais [\textit{over-formal}]''\footnote{Tradução de \textcite[p. 64; destaques meus]{deOliveira2023}.}.

Restaurada a ordem e estabelecido algum executivo mais adequado à consecução devida dos negócios, espera-se que ele “coma seus sapos”\footnote{Expressão famosa nas comunidades de empreendedorismo/self-development: “eat your frogs first”: tirar as coisas difíceis logo do caminho.} antes de sentar-se à mesa para fazer negócios “de verdade”, bem sabendo que seus ativos verdes mais trarão retornos \textit{reputacionais} do que financeiros. Tiradas do caminho as responsabilidades sociais/burocráticas, partem para os negócios propriamente ditos: a alocação ótima de seus portfólios, onde vale tudo, desde tabaco até \textit{health care}, desde armas até \textit{death care}. Não é à toa, afinal, que fundos verdes são tão recheados de ativos armamentícios, ou que mineradoras tão facilmente jactem-se como sustentáveis.\footnote{\textbf{COLOCAR CITAÇÕES DE TELÉSFORO.}. “O Rio? É doce // A Vale? Amarga. // Ai, antes fosse // mais leve a carga.” (Lira Itabirana, Carlos Drummond de Andrade)}

\begin{quote}
    ``O investimento sustentável é uma área confusa das finanças, que muitas vezes significa coisas diferentes para pessoas diferentes. Na maioria das vezes, significa construir carteiras de investimento que excluem categorias questionáveis, como o ``desinvestimento'' em produtores de combustíveis fósseis, numa aparente tentativa de combater as mudanças climáticas. Infelizmente, há uma diferença entre se eximir de algo em que você não deseja participar e lutar ativamente contra algo que você acredita que precisa parar para o bem de todos.'' \cite[p. 9]{Fancy2021}
\end{quote}

E, no fim das contas, ``tanto estardalhaço para nada'': manchetes sobre ``o investidor anti-clima'' bombaram, obviamente colocando-o como o errado da história, enquanto aqueles que estão no meio financeiro \textit{sabem} que ele está certo (de uma forma diferente daquela das manchetes), \textit{sabem} da ``contradição em termos'' da expressão ``finanças sustentáveis'', \textit{sabem} que ele diz o que lhes é evidente todos os dias em suas rotinas de trabalho.
\begin{quote}
    ``...a realidade é que grande parte do que importa para a sociedade simplesmente não afeta o retorno de uma estratégia de investimento particular. Muitas vezes, isso se deve ao horizonte temporal do investimento subjacente: muitas estratégias têm um prazo muito curto, o que significa que as questões ESG de longo prazo não são particularmente relevantes. Às vezes, existem estratégias de dívida que possuem vários seguros e proteções [\textit{insurance and protections}] e que, de qualquer forma, não se preocupam com um período posterior à data de vencimento da dívida. Na maioria das vezes, investimentos caros e de longo prazo relacionados à sustentabilidade são incertos e demoram para gerar potencial de lucro, se é que existe algum.'' \cite[p. 15]{Fancy2021}
\end{quote}

\begin{quote}
    ``As pessoas que tomam essas decisões de alocação de capital são treinadas para reagir rapidamente às mudanças nos rendimentos esperados e na rentabilidade das oportunidades de investimento. Se você dissesse a um gestor de portfólio para reduzir a pegada de carbono de seu portfólio, a maioria concordaria com a cabeça e diria que é importante, mas depois voltaria a alocar capital para os empreendimentos mais rentáveis --- exatamente como são legalmente obrigados e financeiramente incentivados a fazer. Se, por outro lado, uma externalidade negativa for ``internalizada'' pelo governo por meio de, digamos, um imposto sobre a poluição, reduzindo assim a rentabilidade dos grandes emissores de gases de efeito estufa, os alocadores de capital \textit{automaticamente} reagirão o mais rápido possível, alocando menos capital para essas oportunidades agora menos rentáveis.

    ``A maioria dos gestores de portfólio sabe intuitivamente que tudo isso acontece. Um gerente de projetos da BlackRock foi explícito comigo: “Eu acredito nas mudanças climáticas. Se tivéssemos um preço para o carbono, eu reduziria minha pegada de carbono da noite para o dia --- e todos os outros também. Mas não faz sentido fazer isso sozinho e me colocar em desvantagem, e não é o que eu sou legalmente obrigado a fazer ou pelo que sou pago.”'' \cite[p. 24]{Fancy2021}
\end{quote}


\begin{quote}
    ``Muitos dos funcionários de nível operacional [\textit{rank and file employees}] em empresas que geram lucro de maneiras que prejudicam o público não se dão conta disso. Frequentemente, isso ocorre porque trabalham em silos [i.e. isoladamente]. Ou porque o problema é complexo. E de longo prazo. Às vezes, tudo isso junto. Para a alta administração [\textit{senior management}], pode ser tentador acreditar em uma fantasia que se alinha perfeitamente aos seus interesses financeiros. Como disse, certa vez, o escritor americano Upton Sinclair: "É difícil fazer um homem entender algo quando seu salário depende de ele não entendê-lo".'' \cite[p. 26]{Fancy2021}
\end{quote}








\section{A insustentável adaptabilidade do homem}
A consideração que o capital concede à crise climática não tem como ser senão passiva: tudo que lhe foge ao alcance financeiro torna-se mera pré-condição de seus \textit{valuations}, e tudo que foge do horizonte de maturidade de seus ativos não lhes diz respeito em absoluto --- \textit{après moi, le déluge}\footnote{``Depois de mim, o dilúvio'', expressão atribuída ao rei francês Luís XV.}! O capital não se importa com estimativas soltas de como o mundo estará em 2100, e sim em como ele estará daqui a alguns poucos anos. (Ironicamente, a janela temporal de sua paciência em reaver seus investimentos tem se estreitado ao longo das décadas justamente devido à crise climática.)

E não só quanto à faceta ambiental. A questão social (e/ou socioambiental) também é vista como se estivesse detrás de uma vitrine de loja: distante ao toque, mas não da carteira. É uma pena que o mundo gire em função do trabalho escravo no Sudeste Asiático, mas, já que já estão sendo empregados, \textit{por que não desfrutarmos de seus produtos}? A possibilidade de mudança de condições de opressão que, ``pelo bem ou pelo mal'', também são úteis e rentáveis ao resto do mundo, é recebida não só com cansaço — ``\textit{there is no alternative}'', ``\textit{so it goes}'' etc —, como também, e principalmente, de forma reativa, na defensiva — ``então você é contra o progresso/desenvolvimento?'', ``então devemos voltar a ser pobres?'', etc.\footnote{Sempre cai como uma luva, nessas reações defensivas, alguma menção do, entre aspas, ``comunismo'' como um espantalho de como o mundo ``poderia, mas não deveria, ser'': escasso, pobre, estático, ditatorial, etc. Realmente, \textit{quão diferente} é o mundo em que vivemos hoje!}

A questão do racismo ambiental sempre está à espreita das discussões de ESG, e, mesmo quando é explicitamente reconhecida, dificilmente é tomada a sério. É claro: porque essa questão é tacitamente pressuposta por tantos \textit{business as usual}. Um exemplo que o atesta claramente pode ser encontrado no trabalho antropológico de Luisa Reis-Castro, que fez um trabalho de campo em Juazeiro (Bahia) relacionado à liberação de mosquitos geneticamente modificados, uma estratégia de erradicação do \textit{Aedes aegypti}. Essa causa ``nobre'', que facilmente pode ser vista como ``cumprindo princípios ESG'', é mais complexa do que parece:
\begin{quote}
    ``Assim, embora se promova como uma estratégia muito inovadora, a técnica de insetos transgênicos RIDL [sigla para ``Liberação de Insetos Portadores de um gene letal Dominante''] segue uma lógica profundamente enraizada que se concentra no mosquito, em vez de analisar e melhorar as condições sociais, saúde e atendimentos médicos. Isso fica evidente no caso brasileiro. A cidade de Juazeiro, no estado da Bahia --- local escolhido para experimentar essa tecnologia de ponta --- \textit{não possui abastecimento de água encanada} [!].'' \cite[p. 124; grifo meu]{Reis-Castro2013}
\end{quote}

Mosquitos esses que são criados por uma empresa, com o propósito de serem comercializados... ou seja, são mercadorias! Mercadorias que são \textit{testadas} no Brasil, mas produzidas por uma empresa... britânica (Oxitec, um braço comercial da Universidade de Oxford). O Brasil serve neste caso, portanto, como mero \textit{laboratório}, e o racismo ambiental dificilmente oculta-se quando esses experimentos, que se tratam de insetos \textit{que proliferam na água}, são feitos numa região historicamente assolada por secas e econômica/politicamente marginalizada. 

E, não bastasse isso, esses mosquitos tinham uma vida melhor que os trabalhadores locais que os ``produziam'':
\begin{quote}
    ``Para manter a produção padronizada, os mosquitos geneticamente modificados receberam a mesma ração para peixes usada no laboratório inglês onde foram [originalmente] desenvolvidos; porém, no Brasil, tratava-se de uma marca importada e cara. Além disso, essa ração precisava ser moída duas vezes e peneirada para se transformar em um pó bem fino. Essas etapas garantiam que não se formassem torrões, permitindo uma dosagem precisa e uma dissolução mais fácil na água. Enquanto os trabalhadores envolvidos realizavam esse processo árduo e trabalhoso, Jonatan [pseudônimo de um dos trabalhadores locais nesse laboratório] comentou: “Toda essa ração importada e todo esse esforço para alimentá-los”. Após uma breve pausa reflexiva, balançando a cabeça, ele disse: “Esses mosquitos têm uma vida melhor que a minha!”.''\cite[p. 331]{Reis-Castro2021}
\end{quote}

Dentre os maiores pecados capitais na catequese do capitalista está o desperdício de capital, independente de seu funcionamento causar disrupções socioambientais. Por sua própria definição, capital consiste em meios pelos quais a produção futura de riquezas será maior que sua produção vigente no presente\footnote{\textbf{A visão vulgar de capital, ao menos...}}; dessa forma, a destruição de capital não só deixa de efetivar essa produção futura, quanto destrói ``riqueza'' presente — é, portanto, ``um desperdício''! Mais que isso, pode-se concluir: porquanto é uma destruição de riqueza já existente, é um ``crime''! Afinal, para que consertar o que ``não está quebrado'', i.e. o que \textit{ainda funciona}? Não é difícil entender por que os apelos às ``iniciativas próprias'' que os princípios ESG pedem ao capital foram, no máximo, inócuas: se o dano já está feito, então de que vale comportar-se bem? Se já temos tantas plataformas de petróleo montadas, funcionando e trazendo lucros e dividendos e \textit{royalties}, para que diabos vamos desmontá-las? Nesses casos, os óbvios interesses financeiros podem até mesmo mascarar-se de moralidade: para que vamos diminuir o ``bem-estar'' do resto do mundo?\footnote{``Por exemplo, uma CEO pode decidir reduzir a pegada de carbono de sua empresa, mas não pode fazê-lo simplesmente por ser ``justo'' ou ``a coisa certa a se fazer''; a decisão precisa ser justificada em termos dos interesses dos acionistas.'' \cite[p. 16]{Fancy2021}}

Não é à toa, portanto, que o discurso de \textit{adaptabilidade} --- e sinônimos próximos, como robustez e resiliência --- tem tanta prevalência no vocabulário capitalista: onde há riscos, há oportunidade; a necessidade é a mãe da invenção, e, por sorte, também é a mãe dos lucros. Porém, o discurso de ``adaptabilidade'' apresenta-se de duas formas sutilmente distintas: adaptabilidade do capital já montado para novas funções, ou seja, de um capital ainda em funcionamento (logo, um uso mais arraigado deste capital); e adaptabilidade do mercado em incorporar novas atividades e possibilidades (logo, instalação de capital \textit{novo}).




\section{O ESG como \textit{greenwashing} sistêmico}
Por tratar-se de algo oposto ao seu próprio pôr teleológico, o ESG é uma
versão meramente formal de sustentabilidade: o capital restringe suas
ações ``sustentáveis'' àquilo que lhe convém, por exemplo, quando se tratam de
\textit{custos} associados a legislações vigentes (e também ponderando
qual o custo-benefício de evadi-las), em cujo caso ele ajusta seu \textit{budget} tendo em vista tais custos. Quando segue tais ações ``sustentáveis'' de bom grado, ou o faz porque elas lhe trarão algum retorno adequado --- por exemplo, em \textit{visibilidade} (propaganda) ---, ou o faz contra seu \textit{télos} --- em cujo caso não está agindo enquanto capital.

Quando o capital --- falando pela voz de algum de seus apóstolos capitalistas --- fala sobre ``ser sustentável'', sobre ``refrear o consumismo exagerado'', ``descarbonizar o planeta'', ``transição energética'' etc., ele fala como que em duas frequências superpostas: transmite uma mensagem de ``sustentabilidade'' à sociedade civil, e também assegura o \textit{business as usual} àqueles que falam sua língua corporativa e conhecem tais meandros semânticos. Sabem que a finitude do planeta os enclausura, mas não podem dar-se ao luxo de realmente honrar suas promessas, devido à inexorável maré da concorrência\footnote{``Concorrência'' no sentido marxiano: \textbf{Citação de Marx sobre Concorrência. Talvez de Lukács em Reificação}} que os impele adiante. Sabem que a genuína sustentabilidade lhes é fundamentalmente oposta: porque ela limita seu caráter desmedido, porque ela constrange sua expansividade, porque ela refreia sua aceleração. Neste sentido, portanto, o ESG não pode ser senão essencialmente \textit{greenwashing}, um cheque com uma assinatura elegante, porém em branco ademais.

A reação defensiva a essa conclusão é que ``o capital também é capaz de fazer o bem pelo mundo'', afinal, ``\textit{nunca na história da humanidade} a mortalidade infantil foi tão baixa'' e estatísticas do tipo. Trata-se, no fundo, do caráter duplo da mercadoria: embora ela somente possa devir valor porquanto é valor de uso --- somente se troca algo no mercado porquanto essa coisa satisfaça necessidades humanas (as exceções não invalidam esta regra) ---, o capital não a produz \textit{porque} ela é valor de uso, e sim para que ela devenha valor, e, em particular, mais-valor.\footnote{``Aqui {[}no processo capitalista de produção em geral{]}, os
valores de uso só são produzidos porque e na medida em que são o substrato material, os
suportes do valor de troca. E, para nosso capitalista,
trata-se de duas coisas. Primeiramente, ele quer produzir um valor de
uso que tenha um valor de troca [...]. Em segundo lugar,
quer produzir uma mercadoria cujo valor seja
maior do que a soma do valor das mercadorias requeridas para sua produção,
os meios de produção e a força de trabalho, para cuja compra ele
adiantou seu dinheiro no mercado. \textit{Ele quer produzir não só um valor de uso, mas uma mercadoria; não só valor de uso, mas valor, e não só valor, mas também
mais-valor.}''  \cite[p.~263; grifo meu]{Marx2017}} O substrato material sobre o qual assenta-se o valor que produz torna-se quase contingente ao seu objetivo, torna-se até um mal necessário em que tem de incorrer. É irônico, portanto, ver a dissonância entre o \textit{trabalho} corporativo, quantitativo e desapaixonado pelo propósito da própria empresa como ele é por natureza, e o \textit{discurso} corporativo sobre sustentabilidade --- o que o corporativo é dentro de seu escritório \textit{versus} a fachada que ele coloca para os transeuntes olhando de fora. Ou seja, é como se ``trazer valor aos clientes'' fosse \textit{mesmo} aquilo que faz o padeiro de Smith continuar fazendo pães!\footnote{Esse próprio exemplo de Adam Smith gera muita confusão, porque sua analogia faz mais sentido se falar deste padeiro \textit{enquanto produtor capitalista de mercadorias}, não dele enquanto mero indivíduo produzindo bens que vende mais caro do que seus custos.}

É, portanto, ideológico dizer que o capital ``produz em prol do ser humano'': embora ele necessariamente produza valores de uso para clientes que possuam necessidades humanas a serem satisfeitas com seus produtos, ele produz em prol de sua própria valorização, ou seja, em prol de sua própria reprodução. O capital, portanto, assim como Frank Underwood, ``reza \textit{para} si mesmo, \textit{por} si mesmo''\footnote{``\textit{I pray} to \textit{myself}, for \textit{myself}'' (\textit{House of Cards}, T1E12).}. O fato de que o mundo ``é menos pobre hoje do que jamais foi'', e por aí vai, não exclui o fato de que o ímpeto imanente do capital é o de autovalorizar-se --- \textit{este} é seu \textit{télos}. Só porque ele trouxe genuínas comodidades à vida moderna não o torna santo, e não quer dizer que ele produziu tais mercadorias \textit{porque} elas são cômodas: ele as produziu não por bondade, e sim por \textit{ganho}; produziu-as porque tinha quem as comprasse.\footnote{Aqui sim vale a citação original de Smith: ``Não é da benevolência do açougueiro, do cervejeiro ou do padeiro que esperamos nosso jantar, mas da consideração que eles tem pelos próprios interesses. Apelamos não à sua humanidade, mas ao seu amor-próprio, e nunca falamos a eles das nossas próprias necessidades, mas das vantagens que eles podem obter.'' (\textit{Riqueza das Nações}, Cap. II, Livro I)} (Aliás, neste sentido, a \textit{trickle-down economics} acaba sendo mais sincera sobre os motivos próprios do capital do que o \textit{welfare state}.) É irônico, também, que o discurso sobre a crise climática seja tomado como ``alarmista'', ``catastrofista'' e/ou ``moralista'' pelos arautos e \textit{think tanks} do capital, mas que o argumento de que ``o capitalismo é o melhor sistema que temos'', ``o agronegócio abastece o mundo inteiro'', ``democratização do sistema financeiro'' etc. --- as velhas falácias de ``pense nas criancinhas!'' --- seja acatado tão prontamente como motivação para a manutenção do estado atual de coisas.

Tais questionamentos ``tão óbvios'' têm de ser também questionados. ``O consumo energético mundial está cada vez mais crescendo'' --- o que/quem está induzindo esta tendência? É realmente \textit{somente} o crescimento populacional e a democratização da energia elétrica? ``A descarbonização terá disrupções na economia mundial, e não podemos nos dar a este luxo'' --- quais setores serão os mais afetados? Realmente podemos ``nos dar ao luxo'' de postergar a crise inevitável por conta destes setores serem afetados? No fim das contas, quando estas desculpas vêm do capital e de seus seguidores, sempre são ladainhas, sempre são falácias, ou meias-verdades --- e estas últimas \textit{são} as melhores mentiras. Para quem fala sobre o fim do agronegócio: ``quem alimentará um terço do mundo inteiro?''. Para quem fala do fim dos combustíveis fósseis: ``e os países pobres que precisam crescer suas pobres economias?''.  São meias-verdades cobertas por um verniz retórico: as tais das ``pobres economias'' que tanto aparecem nos discursos emocionais do capital não servem para mais do que isso, e para fornecer-lhes matérias-primas baratas, mão-de-obra livre e submissão econômica.\footnote{O governo dos Estados Unidos, sob Trump 2.0, enuncia o que quer ``no mundo, e do mundo'' em seu ``corolário'' à Doutrina Monroe, publicada em novembro de 2025. Dentre outras coisas, ``queremos um Hemisfério [Ocidental] cujos governos cooperem conosco contra narcoterroristas, cartéis e outras organizações criminosas transnacionais; queremos um Hemisfério que permaneça livre de incursões estrangeiras hostis [de quem?] ou da apropriação de ativos essenciais [por quem?], e que apoie cadeias de suprimentos críticas; e queremos garantir nosso acesso contínuo a locais estratégicos importantes [é claro]. Em outras palavras, afirmaremos e aplicaremos um “Corolário Trump” à Doutrina Monroe [sic].'' \cite{WhiteHouse2025a}. Este corolário também foi chamado na mídia, adoravelmente, de ``Doutrina Donroe'', com ``D'' de \textit{yours truly}, Donald Trump...}


Dessa forma, pode-se dizer que o ESG como um todo é um esquema de
\textit{greenwashing}\footnote{``Diversos estudos científicos têm dado razão {[}à percepção de que{]} os fundos ESG, na realidade, não têm impactos ambientais positivos. Não se trata, portanto, de ``separar o joio do trigo'', como se práticas de
\textit{greenwashing} fossem exceções pontuais. O ESG é um discurso e um projeto de \textit{greenwashing} sistêmico, de construção de uma imagem ``verde'' para o capitalismo global, sobretudo para o mercado financeiro, sem promover de modo efetivo a descarbonização da economia, muito menos uma transição justa.'' \cite[p.~80]{Telesforo2025}}, não porque sua aparência dá a entender que sua essência é essa, mas porque sua própria essência é substancialmente ``esquizofrênica'', contraditória. Quando se fala \textit{genuinamente} de sustentabilidade, não fala-se em nome do capital, fala-se de outras coisas; não fala-se dos retornos que advirão de tais e tais projetos, e sim de sua utilidade, de sua importância a comunidades locais etc.; não fala-se de (mais-)valor, e sim de valores de uso. É por isso que é tão comum ver falas do capitalismo verde que comecem com um tom ``militante'' e acabem com o jargão de \textit{business}. Por exemplo:

\begin{quote}
    ``Nosso \textit{Investment Stewardship group} se envolve extensivamente com empresas ao redor do mundo em \textit{questões que são relevantes para a sustentabilidade financeira de longo prazo das empresas} [ou seja, sustentabilidade... \textit{financeira}! É um jogo de palavras comum em discursos desse meio]. Nos últimos anos, escrevemos cartas para CEOs de empresas enfatizando a importância de uma abordagem de longo prazo. Pedimos a eles que articulem suas estratégias de crescimento de longo prazo; \textit{que garantam uma governança adequada; e abordem outras questões sociais e ambientais relevantes a seus modelos de negócios} {[}i.e. é um discurso com tom ESG{]} --- todas estas são etapas que acreditamos indicarem a boa gestão que impulsiona o crescimento de longo prazo {[}se não houver crescimento, \textit{why bother?}{]}. As prioridades da equipe [da Blackrock] para 2018 incluem: estratégia corporativa, governança (incluindo diversidade no \textit{Board}), divulgação de riscos climáticos {[}discurso ESG até agora{]}, \textit{compensation} e políticas de gestão de capital humano {[}discurso de \textit{business}{]}.''(Chairman Letter's de Larry Fink, 2017, p.~20, grifos meus. Disponível em: \url{https://s24.q4cdn.com/856567660/files/oar/2017/downloads/BLK_AR17_Chairmans_Letter.pdf})
\end{quote}

Em 2018, o tom de Larry Fink teve de esfriar: 

\begin{quote}
    ``"É isso que quero dizer quando afirmo que as empresas precisam de modelos \textit{sustentáveis} de negócios {[}grifo no original{]}. \textit{Não se trata de impor valores ambientais ou sociais pessoais às empresas em que investimos em nome dos clientes, nem de a BlackRock de tomar posições políticas} {[}grifo meu; discurso explicitamente ESG{]}. Trata-se de fornecer uma estrutura estratégica e de gestão de riscos que apoie e aprimore a capacidade de uma empresa operar \textit{e gerar valor para seus principais} stakeholders \textit{a longo prazo} [grifo meu; implicitamente quer dizer ``para os \textit{share}holders'', e serve de apaziguamento pelo tom explícito prévio].'' (Chairman Letter's de Larry Fink, 2018, p.~19. Disponível em: \url{https://s24.q4cdn.com/856567660/files/oar/2018/downloads/BlackRock_AR18-CEO-Letter.pdf})
\end{quote}

Em 2021, em sua famosa ``\textit{Letter to CEOs}'' no segundo ano da
pandemia da Covid-19, a questão ESG estava totalmente em alta, e Fink
lança uma convocatória explícita às ``companhias de todos os países'':

\begin{quote}
    ``Embora as questões de raça e etnia variem muito em todo o mundo, esperamos que as empresas em todos os países tenham uma estratégia de talentos que lhes permita aproveitar o conjunto mais amplo de talentos possível. \textbf{Ao publicarem relatórios de sustentabilidade, solicitamos que as divulgações sobre a estratégia de talentos reflitam integralmente os planos de longo prazo para melhorar a diversidade, a equidade e a inclusão, conforme apropriado para cada região} {[}grifo no original!{]}. Nós nos pautamos por esse mesmo padrão {[}!{]}.'' (Disponível em: \url{https://www.blackrock.com/corporate/investor-relations/2021-larry-fink-ceo-letter})
\end{quote}

Será que é tão irônico assim que a Blackrock, logo da segunda vitória eleitoral de Trump no final de 2024, deu com o rabo entre as pernas e silenciosamente mudou
sua página web sobre diversidade, depois, no mínimo, do dia 14 de
dezembro de 2024 \cite{Johnson2025a}? Quão importantes eram suas várias
``\textit{awards and recognitions}'' na área de diversidade e inclusão
(disponíveis nessa página web pré-Trump 2.0), se elas foram tão
incerimoniosamente lançadas ao oblívio?

Como colocado por Baines e Hager,
\begin{quote}
    ``Em suma, adotando uma abordagem abrangente, nossas descobertas indicam que as tentativas mais recentes das \textit{Big Three} [Blackrock, Vanguard, State Street] de distanciarem-se das finanças sustentáveis \textit{são menos uma mudança radical} [\textit{U-turn}] \textit{e mais uma consolidação de uma posição de longa data de obstrucionismo climático} {[}!{]}. Para as \textit{Big Three}, \textit{a gestão ambiental} [\textit{environmental stewardship}] \textit{sempre ficou em segundo plano em relação ao valor para o acionista} [\textit{shareholder value}], \textit{e, portanto, suas recentes tentativas de se distanciarem da retórica das finanças sustentáveis em um contexto de crescente incerteza não são uma surpresa}. Nossa pesquisa, portanto, levanta sérias dúvidas sobre o papel de acionistas como as \textit{Big Three} na defesa do clima [\textit{climate advocacy}]. Na melhor das hipóteses, tais esforços de defesa devem ser considerados um complemento menor a estratégias mais amplas e ambiciosas lideradas pelo Estado para promover uma transição energética de baixo carbono. No entanto, à luz de nossas descobertas, esses esforços provavelmente representam, nas palavras do ex-diretor de investimentos sustentáveis da BlackRock, {[}Tariq Fancy{]}, uma ``distração mortal'' que atrasa tais esforços liderados pelo Estado [...]'' \cite[p.~451-2; grifo meu]{Baines2023}
\end{quote}

E \textit{como} deram para trás: após a vitória eleitoral de Trump em novembro de 2024, já a partir de dezembro houve uma debandada \textit{en masse} de nomes grandes como Goldman Sachs, JP Morgan etc. --- e, claro, Blackrock --- de seus compromissos sustentáveis. 



\section{Considerações finais}
\begin{quote}
    ``Mais do que mero instrumento de propaganda empresarial, o discurso e o aparato institucional ESG foram construídos explicitamente como alternativas a iniciativas regulatórias que impõem custos ao capital ou reduzem o seu poder. Enquanto algumas frações do capitalismo global professam o negacionismo climático puro e simples, outras adotam estratégias diferentes: \textit{reconhecem o problema retoricamente, mas garantem que seu enfrentamento se restrinja às medidas voluntárias do mercado}. O ESG é uma ferramenta de disputa de hegemonia da política climática pelo capital financeiro global, servindo à propagação da racionalidade neoliberal como apta e legítima para dirigir a economia, o Estado e a sociedade.''\cite[p. 82; grifo no original]{Telesforo2025}
\end{quote}

Evidentemente há uma questão de Falácia da Composição aqui: o comprometimento individual não garante que haja um comprometimento coletivo. Assim como dizer que ``a humanidade'' é a principal causadora da atual crise climática não implica que todos os seres humanos tenham a mesma parcela de culpa, a ação meramente individual não muda o fato de que a crise climática é causada social e sistematicamente --- oposto a individual e contingentemente ---, não meramente pela sociedade, mas pela sociedade \textit{capitalista} em particular.

Uma das maiores tragédias do modo de produção capitalista é que a concorrência impele o capital individual à ação mais conveniente para si e detrimental para os demais, ao mesmo tempo que “socializa as perdas” e as responsabilidades pelas ditas “externalidades”; por mais que a crise climática seja causada pelo capitalismo como totalidade, ela certamente é mais atribuível aos grandes conglomerados financeiros e — como mais recentemente tornou-se a nova \textit{auri sacra fames} — a sede por água e a fome por eletricidade dos data centers, do que do estadunidense médio (que não é o mais austero consumidor, diga-se de passagem!).

Para o investidor, o mundo não consiste em mais do que dados e \textit{dashboards} que informam sua alocação de capital; para o capital, são só negócios, alguns mais lucrativos, outros menos; alguns mais voláteis, outros menos. Mas, no fim das contas, todo mundo prefere trabalhar acreditando que está fazendo um bem maior. Afinal, quantos honrados investidores que famintamente abstiveram-se até agora de financiar o genocídio em Gaza suspirarão de alívio quando seu dinheiro puder, \textit{finalmente}, içar suas âncoras e fluir diretamente à “valorosíssima” reconstrução que os Estados Unidos “humildemente” edificarão sobre os escombros da Palestina!\footnote{Trump, em 04/02/2025, disse, em entrevista coletiva, que haveria uma ``oportunidade fenomenal e \textit{unbelievable}'' na Palestina, a ``\textit{Riviera do Oriente Médio}'' --- na qual os palestinos \textit{também} iriam viver \cite{CNBCTelevision2025}. No Fórum Econômico de Davos de 2026, onde Trump anunciou o \textit{Board of Peace} com seus Estados lacaios -- quer dizer, aliados --, Jared Kushner (intermediário nas relações Israel-Gaza e genro de Donald Trump) anunciou o ``plano mestre'' para a reconstrução de Gaza, recortando seu território como se fosse uma planta de urbanismo, com áreas de turismo litorâneo, complexo industrial, parques e áreas residenciais \cite{Haddad2026}. Convenientemente, Kushner é um homem de negócios do setor... imobiliário.} 


O problema do ESG não é sua motivação de tornar o mundo mais sustentável, e sim seu instinto de associar o que ``deve ser feito'' com o que ``traz bons retornos financeiros''. Não é o redirecionamento de recursos financeiros e de capital para atividades mais sustentáveis o que é criticado por uma visão anticapitalista, e sim que estas atividades tenham de fazer sua própria existência caber adequadamente em portfólios de lá e de acolá, que o fardo mais leve da sobrevivência daqueles que delas dependem sejam, antes de tudo, um galardão que mais pese nos bolsos de seus honoráveis patronos. Afinal de contas, poucas coisas no meio corporativo dão o mesmo gozo do que clamar seu merecido galardão sob o clamor de trombetas… mesmo que uma delas eventualmente abra um abismo sob seus pés.\footnote{“E o quinto anjo tocou a sua trombeta, e vi uma estrela que do céu caiu na terra; e foi-lhe dada a chave do poço do abismo. E abriu o poço do abismo, e subiu fumaça do poço, como a fumaça de uma grande fornalha, e com a fumaça do poço escureceu-se o sol e o ar.” (Apocalipse 9:1-2)}

Casos genuínos de sustentabilidade:
\begin{itemize}
    \item Great Green Wall da África
    \item Cidades esponja na China
    \item Billion Oyster Project em Nova York
    \item Movimentos de permacultura, arquitetura sustentável, biomimetismo, etc
    \item Materiais feitos de micélios \cite{Bayer2019}
\end{itemize}


%% \section{A finitude}
%% O capital tem ojeriza à finitude, à restrição. Vê obstáculos não como limites que o impedem, e sim como fronteiras que o instigam sua travessia. Preocupa-se, porém, com o que está logo em sua frente, um problema de cada vez; restringe suas preocupações mediante desconto hiperbólico, e segue sua precificação do \textit{business as usual}. 

%% O que acontece no mundo, o que acomete seus clientes humanos, só diz respeito ao capital porquanto afetar seus negócios: se é positivo a seus clientes, faz mais; se é negativo, se ajusta; e se é oneroso, repensa.


\printbibliography

\end{document}